\documentclass[a4paper, 11pt]{article}

\usepackage[utf8]{inputenc}
\usepackage[T1]{fontenc}
\usepackage[spanish, es-tabla]{babel}
\usepackage{lmodern}
\usepackage[top=2.0cm, bottom=2.0cm, left=2.0cm, right=2.0cm]{geometry}
\usepackage{amsmath, amssymb}
\usepackage{booktabs}
\usepackage{graphicx}
\usepackage{tikz}
\usepackage[dvipsnames]{xcolor}
\usepackage{siunitx}

\sisetup{
    output-decimal-marker = {,},
    separate-uncertainty  = true,
    per-mode              = symbol,
}

\DeclareSIUnit{\division}{div}
\DeclareSIUnit{\ppm}{ppm}
\DeclareSIUnit{\voltpp}{V_{pp}}
\DeclareSIUnit{\voltrms}{V_{rms}}

\definecolor{uni-blue}{RGB}{0, 56, 101}

\pagestyle{empty}

\begin{document}

\begin{center}
    {\Large\bfseries\color{uni-blue} FORMATO OFICIAL PARA RECORTE O COPIA}\\[1.5mm]
    {\normalsize\textbf{Sección 5.8: Tabla Resumen de Especificaciones Metrológicas}}\\[1mm]
    {\footnotesize\textit{Física III (CF1B2) --- Informe de Laboratorio N.° 1 (El Osciloscopio)}}\\[3mm]
    \rule{0.9\linewidth}{0.8pt}
\end{center}

\vspace{0.8cm}

%% CAJA DE RECORTE CON LA TABLA EXACTA
\begin{center}
\begin{tikzpicture}
    \node[draw=gray!70, dashed, line width=0.6pt, inner sep=4mm, fill=white] {
        \begin{minipage}{16.0cm}
            \textbf{\large 5.8. Tabla Resumen de Especificaciones de los Instrumentos}\\[0.25cm]
            {\small En la Tabla~1 se sintetizan las especificaciones metrológicas fundamentales, rangos de operación e incertidumbres instrumentales declaradas para todos los equipos empleados en la experiencia.}\\[0.35cm]
            
            \begin{center}
                \textbf{\normalsize Tabla 1:} \textit{\normalsize Especificaciones metrológicas de los equipos e instrumentos empleados.}
            \end{center}
            \vspace{-0.2cm}
            
            \resizebox{\linewidth}{!}{%
            \begin{tabular}{@{}llll@{}}
                \toprule
                \textbf{Instrumento} & \textbf{Modelo / Marca} & \textbf{Rango de Medida} & \textbf{Incertidumbre Instrumental ($\delta x$)} \\
                \midrule
                Osciloscopio Digital & R\&S HMO1202 / RTC1002 & $\qty{0}{\mega\hertz}\text{ a }\qty{200}{\mega\hertz},\; \qty{1}{\milli\volt\per\division}\text{ a }\qty{10}{\volt\per\division}$ & $\delta N_{\text{div}} = \pm 0{,}1\,\text{div},\; \delta f/f = \pm\qty{0,05}{\percent}$ \\[0.15cm]
                Generador de Señales & OWON AG 051F & $\qty{1}{\micro\hertz}\text{ a }\qty{5}{\mega\hertz},\; \qtyrange{0,02}{20}{\voltpp}$ & $\delta f/f = \pm\qty{1}{\ppm},\; \delta V_p = \pm\qty{0,05}{\volt}$ \\[0.15cm]
                Multímetro Digital   & Industrial True-RMS    & $\qtyrange{0}{600}{\volt}\; (\text{DC/AC})$ & $\delta V_{\text{DC}} = \pm\qty{0,001}{\volt},\; \delta V_{\text{AC}} = \pm\qty{0,01}{\volt}$ \\[0.15cm]
                Transformador AC     & Reductor Monofásico    & $\qty{220}{\volt} \to \qty{6}{\volt},\; \qty{60}{\hertz}$ & $\delta f_{\text{red}} = \pm\qty{0,05}{\hertz}$ \\[0.15cm]
                Fuente Regulada DC   & Multisalida de Banco   & $\qtyrange{2,5}{10,5}{\volt}$ & $\delta V_{\text{mult}} = \pm\qty{0,001}{\volt}$ \\[0.15cm]
                Pila Alcalina        & Cilíndrica AA          & \qty{1,50}{\volt} nominal & $\delta V_{\text{mult}} = \pm\qty{0,001}{\volt}$ \\[0.15cm]
                Cables y Sondas      & Coaxial BNC / Banana   & $\qty{50}{\ohm},\; \qty{600}{\voltrms}\text{ máx.}$ & Atenuación despreciable en baja frecuencia \\
                \bottomrule
            \end{tabular}%
            }
        \end{minipage}
    };
\end{tikzpicture}
\end{center}

\vspace{1.0cm}

\begin{center}
    {\footnotesize\color{gray} \textit{Instrucciones: Recortar por la línea punteada para pegar en la Hoja 10 (Sección 5.8) o fotocopiar directamente.}}
\end{center}

\end{document}
